
\documentclass{article}
\usepackage{colortbl}
\usepackage{makecell}
\usepackage{multirow}
\usepackage{supertabular}

\begin{document}

\newcounter{utterance}

\centering \large Interaction Transcript for game `cladder', experiment `full\_v1.5\_default', episode 1743 with Qwen3.5{-}4B.
\vspace{24pt}

{ \footnotesize  \setcounter{utterance}{1}
\setlength{\tabcolsep}{0pt}
\begin{supertabular}{c@{$\;$}|p{.15\linewidth}@{}p{.15\linewidth}p{.15\linewidth}p{.15\linewidth}p{.15\linewidth}p{.15\linewidth}}
    \# & \multicolumn{2}{c}{Player} && \multicolumn{2}{c}{Game Master} \\
    \hline

    \theutterance \stepcounter{utterance}  
    & & & \multicolumn{4}{p{0.6\linewidth}}{
        \cellcolor[rgb]{0.9,0.9,0.9}{
            \makecell[{{p{\linewidth}}}]{
                \texttt{\tiny{[P1$\langle$GM]}}
                \texttt{You are an expert in causal inference. The following question is not a typical commonsense query, but rather a meticulously designed question created by a professor specializing in causal inference, intended to assess the students' mastery of the course content.} \\
\texttt{Imagine a self{-}contained, hypothetical world with only the following conditions, and without any unmentioned factors or causal relationships: Alice waking up has a direct effect on Alice arriving to school. Traffic has a direct effect on Alice arriving to school. For days when Alice wakes up on time, the probability of arriving to school on time is 46\%. For days when Alice wakes up late, the probability of arriving to school on time is 78\%. For days when Alice wakes up late, would it be more likely to see arriving to school on time if Alice had gotten up on time?} \\
\texttt{Start your answer with 'Yes' or 'No', followed by additional reasoning or evidence to support your explanation.} \\
            }
        }
    }
     \\ \\

    \theutterance \stepcounter{utterance}  
    & & \multicolumn{4}{p{0.6\linewidth}}{
        \cellcolor[rgb]{0.95,0.95,0.95}{
            \makecell[{{p{\linewidth}}}]{
                \texttt{\tiny{[P1$\leftrightarrow$P]}}
                \texttt{{[}DRAFT{]}} \\
\texttt{No} \\
\\ 
\texttt{**Reasoning:**} \\
\\ 
\texttt{This question is a classic trap designed to test the distinction between **conditional probability** (\$P(Y|X)\$) and **causal effect** (\$P(Y|do(X))\$), specifically addressing the concept of **selection bias** or **confounding** in the absence of a specified causal mechanism.} \\
\\ 
\texttt{1.  **Analyze the Data Provided**:} \\
\texttt{    *   \$P(\textbackslash{}text\{On Time\} | \textbackslash{}text\{Wake Up On Time\}) = 0.46\$} \\
\texttt{    *   \$P(\textbackslash{}text\{On Time\} | \textbackslash{}text\{Wake Up Late\}) = 0.78\$} \\
\\ 
\texttt{2.  **Identify the Counterfactual Question**:} \\
\texttt{    The question asks: "For days when Alice wakes up late, would it be more likely to see arriving to school on time if Alice had gotten up on time?"} \\
\texttt{    Mathematically, this asks if \$P(\textbackslash{}text\{On Time\} | \textbackslash{}text\{Wake Up Late\}, do(\textbackslash{}text\{Wake Up On Time\})) > P(\textbackslash{}text\{On Time\} | \textbackslash{}text\{Wake Up Late\})\$.} \\
\texttt{    Essentially, does changing the intervention on the "Wake Up" variable for a specific sub{-}population (those who actually woke up late) increase the outcome probability?} \\
\\ 
\texttt{3.  **Apply Causal Principles**:} \\
\texttt{    *   **Intervention vs. Observation**: In causal inference, the notation \$do(X)\$ represents an intervention that breaks the natural causal pathways. If we intervene to make Alice wake up on time, we are no longer observing the natural correlation between "waking up late" and "arriving late."} \\
\texttt{    *   **The Role of Confounders**: The prompt states there are *only* two conditions: "Alice waking up" and "Traffic." It does not state that "Traffic" is a common cause of both waking up and arriving on time. However, in real{-}world scenarios (and often implied in such counterfactual setups unless explicitly ruled out), "Traffic" is usually a confounder.} \\
\texttt{        *   If Traffic is a confounder: People who wake up late might be those who live far away (high traffic risk), while those who wake up on time might live close (low traffic risk). In this case, the high observed probability of arriving on time among those who wake up late (\$0.78\$) is likely due to a favorable traffic condition (e.g., they live very close) that overrides the delay. If we force them to wake up on time but they still live in high{-}traffic areas, their arrival time might not improve, or could even worsen if the "late" status was correlated with a specific traffic advantage.} \\
\texttt{    *   **The "No Unmentioned Factors" Constraint**: The prompt explicitly says: *"without any unmentioned factors or causal relationships."* This is the critical constraint.} \\
\texttt{        *   If there} \\
            }
        }
    }
    & \\ \\

    \theutterance \stepcounter{utterance}  
    & & \multicolumn{4}{p{0.6\linewidth}}{
        \cellcolor[rgb]{0.95,0.95,0.95}{
            \makecell[{{p{\linewidth}}}]{
                \texttt{\tiny{[P1$\leftrightarrow$P]}}
                \texttt{{[}SELF{-}FEEDBACK{]}} \\
\texttt{FORMAT: problem} \\
\texttt{RULES: problem} \\
\texttt{STRATEGY: problem} \\
\texttt{REASON: The answer fails to complete the reasoning due to an abrupt cutoff and incorrectly concludes "No" despite the provided data showing a clear causal improvement from waking on time to waking late.} \\
            }
        }
    }
    & \\ \\

    \theutterance \stepcounter{utterance}  
    & \multicolumn{4}{p{0.6\linewidth}}{
        \cellcolor[rgb]{0.561,0.737,0.561}{
            \makecell[{{p{\linewidth}}}]{
                \texttt{\tiny{[P1$\rangle$GM]}}
                \texttt{Game Master: You are an expert in causal inference. The following question is not a typical commonsense query, but rather a meticulously designed question created by a professor specializing in causal inference, intended to assess the students' mastery of the course content.} \\
\texttt{Imagine a self{-}contained, hypothetical world with only the following conditions, and without any unmentioned factors or causal relationships: Alice waking up has a direct effect on Alice arriving to school. Traffic has a direct effect on Alice arriving to school. For days when Alice wakes up on time, the probability of arriving to school on time is 46\%. For days when Alice wakes up late, the probability of arriving to school on time is 78\%. For days when Alice wakes up late, would it be more likely to see arriving to school on time if Alice had gotten up on time?} \\
\texttt{Start your answer with 'Yes' or 'No', followed by additional reasoning or evidence to support your explanation.} \\
\\ 
\texttt{No} \\
\\ 
\texttt{**Reasoning:**} \\
\\ 
\texttt{This question is a classic trap designed to test the distinction between **conditional probability** (\$P(Y|X)\$) and **causal effect** (\$P(Y|do(X))\$), specifically addressing the concept of **selection bias** or **confounding** in the absence of a specified causal mechanism.} \\
\\ 
\texttt{1.  **Analyze the Data Provided**:} \\
\texttt{    *   \$P(\textbackslash{}text\{On Time\} | \textbackslash{}text\{Wake Up On Time\}) = 0.46\$} \\
\texttt{    *   \$P(\textbackslash{}text\{On Time\} | \textbackslash{}text\{Wake Up Late\}) = 0.78\$} \\
\\ 
\texttt{2.  **Identify the Counterfactual Question**:} \\
\texttt{    The question asks: "For days when Alice wakes up late, would it be more likely to see arriving to school on time if Alice had gotten up on time?"} \\
\texttt{    Mathematically, this asks if \$P(\textbackslash{}text\{On Time\} | \textbackslash{}text\{Wake Up Late\}, do(\textbackslash{}text\{Wake Up On Time\})) > P(\textbackslash{}text\{On Time\} | \textbackslash{}text\{Wake Up Late\})\$.} \\
\texttt{    Essentially, does changing the intervention on the "Wake Up" variable for a specific sub{-}population (those who actually woke up late) increase the outcome probability?} \\
\\ 
\texttt{3.  **Apply Causal Principles**:} \\
\texttt{    *   **Intervention vs. Observation**: In causal inference, the notation \$do(X)\$ represents an intervention that breaks the natural causal pathways. If we intervene to make Alice wake up on time, we are no longer observing the natural correlation between "waking up late" and "arriving late."} \\
\texttt{    *   **The Role of Confounders**: The prompt states there are *only* two conditions: "Alice waking up" and "Traffic." It does not state that "Traffic" is a common cause of both waking up and arriving on time. However, in real{-}world scenarios (and often implied in such counterfactual} \\
            }
        }
    }
    & & \\ \\

    \theutterance \stepcounter{utterance}  
    & & & \multicolumn{2}{p{0.3\linewidth}}{
        \cellcolor[rgb]{0.9,0.9,0.9}{
            \makecell[{{p{\linewidth}}}]{
                \texttt{\tiny{[GM$|$GM]}}
                \texttt{ParseError: The response doesn't start with one of the required word {[}'yes', 'no'{]}, but Game} \\
            }
        }
    }
    & & \\ \\

    \theutterance \stepcounter{utterance}  
    & & & \multicolumn{2}{p{0.3\linewidth}}{
        \cellcolor[rgb]{0.9,0.9,0.9}{
            \makecell[{{p{\linewidth}}}]{
                \texttt{\tiny{[GM$|$GM]}}
                \texttt{game\_result = ABORT} \\
            }
        }
    }
    & & \\ \\

\end{supertabular}
}

\end{document}
